\documentclass[12pt,a4paper,reqno]{amsart}

% language
\usepackage[greek,english]{babel}
\usepackage[utf8]{inputenc}
\usepackage{alphabeta}

% change default names to greek
\addto\captionsenglish{
    \renewcommand{\contentsname}{Περιεχόμενα}
    \renewcommand{\refname}{Βιβλιογραφία}
    \renewcommand{\datename}{Ημερομηνία:}
    \renewcommand{\urladdrname}{Ιστοσελίδα}
}

% math 
\usepackage{amsmath,amsthm,amssymb,amscd}

% font
\usepackage[cal=euler]{mathalfa}
\usepackage{libertinus-type1}
% \usepackage{txfonts} % for upright greek letters
\usepackage{bm} % for bold symbols
\usepackage{bbm} % for the simply-looking bb symbols

% miscellaneous 
\usepackage{changepage} %for indenting environments
\usepackage{csquotes} % example: \textcquote{}
\usepackage{blkarray}
\setcounter{MaxMatrixCols}{20} % default for pmatrix is 10!!
\usepackage{ytableau}
\usepackage{array} %needed to increase the vertical length in a tabular

% drawing
\usepackage{tikz,tikz-cd}
\usetikzlibrary{shapes.misc, patterns, matrix, calc, intersections,positioning,arrows.meta}
\usepackage{graphics,graphicx}
\usepackage{float} % provides enhanced control and customization options for floating objects such as figures and tables

% colors
\usepackage{xcolor}
\definecolor{darkcandyapplered}{rgb}{0.64, 0.0, 0.0}
\definecolor{midnightblue}{rgb}{0.1, 0.1, 0.44}
\definecolor{mylightblue}{HTML}{336699}
\definecolor{burntorange}{rgb}{0.8, 0.33, 0.0}
\definecolor{iceberg}{rgb}{0.44, 0.65, 0.82}
\definecolor{applegreen}{rgb}{0.55, 0.71, 0.0}
\definecolor{canaryyellow}{rgb}{1.0, 0.94, 0.0}
\definecolor{lightgray}{rgb}{0.83, 0.83, 0.83}

% hrefs
\usepackage{hyperref}
\usepackage[noabbrev,capitalize]{cleveref}
\hypersetup{
    pdftoolbar=true,        
    pdfmenubar=true,        
    pdffitwindow=false,     
    pdfstartview={FitH},  % fits the width of the page to the window
    pdftitle={},
    pdfauthor={},
    pdfsubject={},
    pdfkeywords={},
    pdfnewwindow=true,  % links in new window
    colorlinks=true,  % false: boxed links; true: colored links
    linkcolor=darkcandyapplered,   % color of internal links
    citecolor=midnightblue,  % color of links to bibliography
    urlcolor=cyan,  % color of external links
    linktocpage=true  % changes the links from the section body to the page number
    }

% geometry
\textwidth=16cm 
\textheight=21cm 
\hoffset=-55pt 
\footskip=25pt

% thm envs (you might need to change the path)
\theoremstyle{definition}
\newtheorem{exercise}{Άσκηση}
\newtheorem{solution}{Λύση}
\newtheorem*{remark}{Παρατήρηση}
\newtheorem*{example}{Παράδειγμα}

% math ops 
% fields
\newcommand{\NN}{\mathbbmss{N}} 
\newcommand{\ZZ}{\mathbbmss{Z}} 
\newcommand{\QQ}{\mathbbmss{Q}} 
\newcommand{\RR}{\mathbbmss{R}} 
\newcommand{\CC}{\mathbbmss{C}} 
\newcommand{\KK}{\mathbbmss{K}} 
\newcommand{\FF}{\mathbbmss{F}} 
\newcommand{\HH}{\mathbbmss{H}} 

% symmetric group
\newcommand{\fS}{\mathfrak{S}}  
\newcommand{\fp}{\mathfrak{p}}  
\newcommand{\fm}{\mathfrak{m}}  

% calligraphic 
\newcommand{\aA}{\mathcal{A}} 
\newcommand{\bB}{\mathcal{B}}
\newcommand{\cC}{\mathcal{C}}
\newcommand{\dD}{\mathcal{D}}
\newcommand{\eE}{\mathcal{E}}
\newcommand{\fF}{\mathcal{F}}
\newcommand{\hH}{\mathcal{H}}
\newcommand{\iI}{\mathcal{I}}
\newcommand{\lL}{\mathcal{L}}
\newcommand{\nN}{\mathcal{N}}
\newcommand{\oO}{\mathcal{O}}
\newcommand{\pP}{\mathcal{P}}
\newcommand{\sS}{\mathcal{S}}
\newcommand{\mM}{\mathcal{M}}
\newcommand{\uU}{\mathcal{U}}

% bold
\newcommand{\bfa}{\mathbf{a}}
\newcommand{\bfe}{\mathbf{e}}
\newcommand{\bfF}{\pmb{F}}
\newcommand{\bfR}{\pmb{R}}
\newcommand{\bfv}{\mathbf{v}}
%\newcommand{\bfx}{\bm{x}}
%\newcommand{\bfx}{\mathbf{x}} 
\newcommand{\bfx}{\pmb{x}}
\newcommand{\bfX}{\pmb{X}}
\newcommand{\bfy}{\pmb{y}}
\newcommand{\bfz}{\pmb{z}}

% roman
\newcommand{\rmA}{\mathrm{A}}
\newcommand{\rmB}{\mathrm{B}}
\newcommand{\rmC}{\mathrm{C}}
\newcommand{\rmD}{\mathrm{D}} 
\newcommand{\rmH}{\mathrm{H}} 
\newcommand{\rmh}{\mathrm{h}} 
\newcommand{\rmI}{\mathrm{I}} 
\newcommand{\rmJ}{\mathrm{J}} 
\newcommand{\rmK}{\mathrm{K}}
\newcommand{\rmM}{\mathrm{M}}
\newcommand{\rmP}{\mathrm{P}}  
\newcommand{\rmp}{\mathrm{p}}  
\newcommand{\rmQ}{\mathrm{Q}}  
\newcommand{\rmR}{\mathrm{R}}
\newcommand{\rmS}{\mathrm{S}}
\newcommand{\rmT}{\mathrm{T}}
\newcommand{\rmU}{\mathrm{U}}
\newcommand{\rmV}{\mathrm{V}}
\newcommand{\rmY}{\mathrm{Y}}
\newcommand{\rmZ}{\mathrm{Z}}
\newcommand{\rmz}{\mathrm{z}}

% arrows and symbols 
\renewcommand{\to}{\rightarrow}
\newcommand{\toto}{\longrightarrow}
\newcommand{\mapstoto}{\longmapsto}
\newcommand{\then}{\Rightarrow}
\newcommand{\IFF}{\Leftrightarrow}
\newcommand{\tl}{\tilde}
\newcommand{\wtl}{\widetilde}
\newcommand{\ol}{\overline}
\newcommand{\ul}{\underline}
\newcommand{\oldemptyset}{\emptyset}
\renewcommand{\emptyset}{\varnothing}
\DeclareMathSymbol{\Arg}{\mathbin}{AMSa}{"39} % for arguments 
\newcommand{\onto}{\ensuremath{\twoheadrightarrow}}
\newcommand{\tle}{\trianglelefteq}
\newcommand{\tge}{\trianglerighteq}

% custom ops
\newcommand{\rmKer}{\mathrm{Ker}}
\newcommand{\rmIm}{\mathrm{Im}}
\newcommand{\lcm}{\mathrm{lcm}}
\newcommand{\GL}{\mathrm{GL}}
\newcommand{\ev}{\mathrm{ev}}
\newcommand{\End}{\mathrm{End}}
\newcommand{\rmid}{\mathrm{id}}
\newcommand{\rmspan}{\mathrm{span}}
\newcommand{\Spec}{\mathrm{Spec}}
\newcommand{\mSpec}{\mathrm{mSpec}}
\newcommand{\rad}{\mathrm{rad}}
\newcommand{\Rad}{\mathrm{Rad}}
\newcommand{\Jac}{\mathrm{Jac}}

% absolute value symbol
\usepackage{mathtools} 
\DeclarePairedDelimiter\abs{\lvert}{\rvert}%
\DeclarePairedDelimiter\norm{\lVert}{\rVert}%
\makeatletter
\let\oldabs\abs
\def\abs{\@ifstar{\oldabs}{\oldabs*}}

% tensor symbol
\newcommand{\tensor}[1]{%
  \mathbin{\mathop{\otimes}\limits_{#1}}%
}

% permutation cycle notation
\ExplSyntaxOn
\NewDocumentCommand{\cycle}{ O{\;} m }
 {
  (
  \alec_cycle:nn { #1 } { #2 }
  )
 }

\seq_new:N \l_alec_cycle_seq
\cs_new_protected:Npn \alec_cycle:nn #1 #2
 {
  \seq_set_split:Nnn \l_alec_cycle_seq { , } { #2 }
  \seq_use:Nn \l_alec_cycle_seq { #1 }
 }
\ExplSyntaxOff

% setminus symbol
\newcommand{\mysetminusD}{\hbox{\tikz{\draw[line width=0.6pt,line cap=round] (3pt,0) -- (0,6pt);}}}
\newcommand{\mysetminusT}{\mysetminusD}
\newcommand{\mysetminusS}{\hbox{\tikz{\draw[line width=0.45pt,line cap=round] (2pt,0) -- (0,4pt);}}}
\newcommand{\mysetminusSS}{\hbox{\tikz{\draw[line width=0.4pt,line cap=round] (1.5pt,0) -- (0,3pt);}}}
\newcommand{\sm}{\mathbin{\mathchoice{\mysetminusD}{\mysetminusT}{\mysetminusS}{\mysetminusSS}}}

% miscellaneous commands
\newcommand{\defn}[1]{{\color{mylightblue}{#1}}}
\newcommand{\toDo}{{\bf\color{red} TODO}}
\newcommand{\toCite}{{\bf\color{green} CITE}}
\newcommand*{\vertbar}{\rule[-1ex]{0.5pt}{2.5ex}} % for matrices with column vectors
\newcommand*{\horzbar}{\rule[.5ex]{2.5ex}{0.5pt}} % for matrices with row vectors
\newcommand{\myblue}[1]{{\color{iceberg}{#1}}}
\newcommand{\myorange}[1]{{\color{burntorange}{#1}}}
\newcommand{\mygreen}[1]{{\color{applegreen}{#1}}}
\newcommand{\myred}[1]{{\color{darkcandyapplered}{#1}}}

% ferrer's diagram
\newcommand{\fdiagram}[1]{
    \begin{tikzpicture}[scale=.7]
        \fill foreach \Z [count=\Y] in {#1}
        {foreach \X in {1,...,\Z} 
        {(\X,-\Y) circle[radius=3pt]}};
    \end{tikzpicture}
}

%
\usepackage{pgfplots}
\pgfplotsset{compat=1.18}

%
\makeatletter
\def\mathcolor#1#{\@mathcolor{#1}}
\def\@mathcolor#1#2#3{%
  \protect\leavevmode
  \begingroup
    \color#1{#2}#3%
  \endgroup
}
\makeatother

% 
\newenvironment{nouppercase}{%
  \let\uppercase\relax%
  \renewcommand{\uppercasenonmath}[1]{}}{}

% titlepage
\title{226: Θεωρία Δακτυλίων και Modules \\ Ασκήσεις με Ενδεικτικές Λύσεις}
\author[Β.~Δ. Μουστακας]{Βασίλης Διονύσης Μουστάκας \\ Πανεπιστήμιο Κρήτης}
\date{11 Μαρτίου 2026}
% \urladdr{\href{https://sites.google.com/view/vasmous}{https://sites.google.com/view/vasmous}}

\begin{document}

\begingroup
\def\uppercasenonmath#1{} % this disables uppercase title
\let\MakeUppercase\relax % this disables uppercase authors
\maketitle
\endgroup

\setcounter{exercise}{10}
% \setcounter{theorem}{5}
\setcounter{equation}{3}
\renewcommand{\thesubsection}{\thesection.\arabic{subsection}}

Σε ότι ακολουθεί υποθέτουμε ότι $R$ είναι ένας \emph{μεταθετικός} δακτύλιος και $n$ είναι ένας θετικός ακέραιος, εκτός αν αναφέρεται διαφορετικά.

\begin{exercise}{\rm(\cite[Άσκηση~7.4.30]{DF04})}
  \label{ex:DF_7.4.30}
  Έστω $I$ ένα ιδεώδες του $R$. Το σύνολο\footnote{Στο \cite{DF04} συμβολίζεται με $\rad(I)$.} 
  \[
  \sqrt{I} \coloneqq \{r \in R : r^n \in I, \ \text{για κάποιο $n \ge 1$}\}
  \]
  ονομάζεται \defn{ριζικό} (radical) του $I$. Να δείξετε τα εξής.
  \begin{itemize}
    \item[(α)] Έχουμε $\sqrt{0} = \nN(R)$, όπου $0$ είναι το τετριμμένο ιδεώδες.
    \item[(β)] Το $\sqrt{I}$ είναι ιδεώδες του $R$ που περιέχει το $I$.
    \item[(γ)] Έχουμε $\sqrt{I}/I = \nN(R/I)$.
  \end{itemize}
  Υπενθυμίζουμε ότι με $\nN(R)$ συμβολίζουμε το μηδενοριζικό του $R$ (βλ. Άσκηση 7).
\end{exercise}

\begin{proof}[Λύση]
  Το (α) έπεται άμεσα από τους ορισμούς και δικαιολογεί τον όρο μηδενοριζικό. Για το (β), αρχικά παρατηρούμε ότι $I \subseteq \sqrt{I}$. Έπειτα, θα επαληθεύσουμε τον Ορισμό 1.5.
  \begin{itemize}
    \item Το $\sqrt{I}$ περιέχει το μηδέν, διότι περιέχει το $I$ που το περιέχει ως ιδεώδες.
    \item Το $\sqrt{I}$ είναι κλειστό ως προς την πρόσθεση. Πράγματι, έστω $a, b \in \sqrt{I}$. Τότε $a^n, b^m \in I$ για κάποια $n, m \ge 1$. Από την Άσκηση 3 (διωνυμικό θεώρημα) έπεται ότι 
    \begin{align*}
      (a+b)^{n+m-1} 
      &= \sum_{k=0}^{n+m-1}\binom{m+n-1}{k} a^kb^{m+n-1-k} \\
      &= \underbrace{b^m\left(
        \sum_{k=0}^{n-1} \binom{n+m-1}{k} a^kb^{n-1-k} 
      \right)}_{\in \ I} \\
      &\quad + \underbrace{a^n\left(
        \sum_{k=n}^{n+m-1} \binom{n+m-1}{k} a^{k-n}b^{n+m-1-k}
      \right)}_{\in \ I}
    \end{align*}
    και γι αυτό $a + b \in \sqrt{I}$, διότι το $I$ είναι ιδεώδες.
    \item Το $\sqrt{I}$ είναι κλειστό ως προς πολλαπλασιασμό με στοιχείο του $R$. Πράγματι, έστω $a \in \sqrt{I}$ και $r \in R$. Τότε $a^n \in I$, για κάποιο $n \ge 1$ και γι αυτό 
    \[
    (ra)^n = r^n a^n \in I
    \]
    διότι το $I$ είναι ιδεώδες.
  \end{itemize}

  Για το (γ), έστω $a + I \in \nN(R/I)$. Τότε
  \[
  0_R + I = (a+I)^n = a^n + I
  \]
  για κάποιο $n \ge 1$. Ισοδύναμα $a^n \in I$, για κάποιο $n \ge 1$ και γι αυτό $a \in \sqrt{I}$ που σημαίνει ότι $a + I \in \sqrt{I}/I$. Όλα τα βήματα αντιστρέφονται και το ζητούμενο έπεται.
\end{proof}

\begin{exercise}{\rm(\cite[Άσκηση~7.4.31]{DF04})}
  \label{ex:DF_7.4.31}
  Έστω $I$ ένα ιδεώδες του $R$. Το $I$ ονομάζεται \defn{ριζικό} αν $\sqrt{I} = I$.
  \begin{itemize}
    \item[(α)] Να δείξετε ότι κάθε πρώτο ιδεώδες του $R$ είναι ριζικό.
    \item[(β)] Να δείξετε ότι το $n\ZZ$ είναι ριζικό στο $\ZZ$ αν και μόνο αν ο $n$ είναι \emph{ελεύθερος τετραγώνων} (square free), δηλαδή αν μπορεί να γραφεί ως γινόμενο διακεκριμένων πρώτων.  
  \end{itemize}
\end{exercise}

\begin{proof}[Λύση]
  Υπενθυμίζουμε ότι με $\Spec(R)$ συμβολίζουμε το σύνολο όλων των πρώτων ιδεωδών του $R$.
  \begin{itemize}
    \item[(α)] Έστω $\fp \in \Spec(R)$. Από την Πρόταση 1.11 έπεται ότι ο $R/\fp$ είναι ακέραια περιοχή και γι αυτό $\nN(R/\fp) = 0$ από την Άσκηση 9. Από την Άσκηση \ref{ex:DF_7.4.30} έπεται ότι 
    \[
    \sqrt{\fp}/\fp = \nN(R/\fp) = 0
    \]
    που σημαίνει ότι $\sqrt{\fp} = \fp$ και κατά συνέπεια το $\fp$ είναι ριζικό.
    \item[(β)] Έστω $n = p_1^{k_1}p_2^{k_2}\cdots p_m^{k_m}$ η παραγοντοποίηση του $n$ σε πρώτους. Συνδυάζοντας τις Ασκήσεις 9 και \ref{ex:DF_7.4.30} προκύπτει ότι 
    \[
    \sqrt{n\ZZ}/n\ZZ = \nN(\ZZ/n\ZZ) = (p_1p_2\cdots p_m)\ZZ/n\ZZ.
    \]
    Συνεπώς,  
    \[
    \sqrt{n\ZZ} = (p_1p_2\cdots p_m)\ZZ
    \]
    και γι αυτό $\sqrt{n\ZZ} = n\ZZ$ αν και μόνο αν $k_1 = k_2 = \cdots = k_m = 1$.
  \end{itemize}
\end{proof}

\begin{exercise}{\rm(\cite[Άσκηση~7.4.32]{DF04})}
  \label{ex:DF_7.4.32}
  Έστω $\mSpec(R)$ το σύνολο όλων των μεγιστικών ιδεωδών του $R$. Το σύνολο 
  \[
  \Jac(I) \coloneqq \bigcap_{\substack{\fm \in \mSpec(R) \\ I \subseteq \fm}} \fm
  \]
  ονομάζεται \defn{ριζικό Jacobson} του $I$. Το ριζικό Jacobson του τετριμμένου ιδεώδους ονομάζεται ριζικό Jacobson του $R$\footnote{Και συχνά στη βιβλιογραφία συμβολίζεται ως $\Rad(R)$ ή $\rmJ(R)$.}.
  \begin{itemize}
    \item[(α)] Να δείξετε ότι το ριζικό Jacobson του $R$ περιέχει το μηδενοριζικό του.
    \item[(γ)] Να υπολογιστεί το $\Jac(n\ZZ)$.
  \end{itemize}
\end{exercise}

\begin{proof}[Λύση]
  Για το (α), θυμόμαστε ότι κάθε μεγιστικό ιδεώδες του $R$ είναι πρώτο, δηλαδή 
  \[
  \mSpec(R) \subseteq \Spec(R).
  \] 
  Αυτό έχει ως συνέπεια 
  \[
  \nN(R) \subseteq \bigcap_{\fp \in \Spec(R)} \fp \subseteq \bigcap_{\fm \in \mSpec(R)} \fm = \Jac(0)
  \]
  όπου ο πρώτος εγκλεισμός έπεται από την Άσκηση 10. 

  Για το (γ), θυμόμαστε ότι 
  \[
  \mSpec(\ZZ) = \Spec(\ZZ) = \{ p\ZZ : \ \text{ο $p$ είναι πρώτος}\}.
  \]
  Άρα, 
  \[
  \Jac(n\ZZ) 
  = \bigcap_{\substack{\text{$p$ πρώτος} \\ n\ZZ \subseteq p\ZZ}} p\ZZ
  = \bigcap_{\substack{\text{$p$ πρώτος} \\ p \mid n}} p\ZZ
  = (p_1p_2\cdots{p_k})\ZZ
  \]
  όπου $p_1, p_2, \dots, p_k$ είναι οι διακεκριμένοι πρώτοι που εμφανίζονται στην παραγοντοποίηση του $n$ σε πρώτους.
\end{proof}

\begin{exercise}{\rm(\cite[Άσκηση~9.2.7]{DF04})}
  \label{ex:DF_9.2.7}
  Να περιγράψετε τα ιδεώδη του $\ZZ[x]/(2, x^3+1)$.
\end{exercise}

\begin{proof}[Λύση]
  Γράφουμε $J = (x^3+1)$. Θεωρούμε την απεικόνιση 
  \begin{align*}
    \ZZ[x] &\to \ZZ_2[x] / (x^3+1) \\
    a_0 + a_1x + \cdots + a_nx^n &\mapsto (\ol{a_0} + \ol{a_1}x + \cdots + \ol{a_n}x^n) + J,
  \end{align*}
  όπου $\ol{a_i} = a_i + 2\ZZ$, για κάθε $i \ge 0$. Αυτή είναι επιμορφισμός δακτυλίων με πυρήνα το ιδεώδες $(2, x^3+1)$ (γιατί;). Συνεπώς, από το πρώτο θεώρημα ισομορφισμού, έπεται ότι\footnote{Ισοδύναμα, μπορεί κανείς να επικαλεστεί το τρίτο θεώρημα ισομορφισμού και να γράψει $\ZZ[x]/(2,x^3+1) \cong \frac{\ZZ[x]/(2)}{(2,x^3+1)/(2)} \cong \ZZ_2[x]/(x^3+1)$.} 
  \[
  \ZZ[x]/(2, x^3+1) \cong \ZZ_2[x] / (x^3+1).
  \]

  Κατά συνέπεια, τα ιδεώδη του $\ZZ[x]/(2, x^3+1)$ αντιστοιχούν στα ιδεώδη του $\ZZ_2[x] / (x^3+1)$, τα οποία με τη σειρά τους αντιστοιχούν στα ιδεώδη $I$ του $\ZZ_2[x]$ τα οποία περιέχουν το $(x^3+1)$, από το τέταρτο θεώρημα ισομορφισμού. Όμως, το $\ZZ_2[x]$ είναι περιοχή κύριων ιδεωδών, από το Θεώρημα 3.3, διότι το $\ZZ_2$ είναι σώμα και γι αυτό $I = (f(x))$ για κάποιο $f(x) \in \ZZ_2[x]$. Συνεπώς, η συνθήκη $(x^3 + 1) \subseteq I = (f(x))$ γίνεται $f(x) \mid x^3 + 1$.

  Η παραγοντοποίηση του $x^3+1$ σε ανάγωγα πολυώνυμα του $\ZZ_2[x]$ είναι 
  \[
  x^3 + 1 = (x+1)(x^2 + x + 1)
  \]
  (γιατί;) και γι αυτό προκύπτουν τέσσερα διαφορετικά ιδεώδη $I$, ένα για κάθε διαίρετη του $x^3+1$. Πιο συγκεκριμένα, από την Πρόταση 3.9, υπολογίζουμε 
  \begin{align*}
  \ZZ_2[x]/(x^3+1) 
  &= \{\ol{a_0}(1+J) + \ol{a_1}(x + J) + \ol{a_2}(x^2 + J) : \ol{a_i} \in \ZZ_2, \ \text{για κάθε $0 \le i \le 2$}\} \\
  &= \{0 + J, 1 + J, x + J, (1 + x) + J, x^2 + J, (x + x^2) + J, (x^2 + x + 1) + J\},
  \end{align*}
  όπου έχουμε ταυτίσει τους συντελεστές $0$ και $1$ με τις εικόνες τους $\ol{0}$ και $\ol{1}$ στο $\ZZ_2$. Κατά συνέπεια, έχουμε 
  \[
    \renewcommand{\arraystretch}{1.3} 
    \begin{tabular}{c|c|c}
    $f(x)$                & $I/J$                                           & γεννήτορας       \\ \hline
    $1$                   & $\ZZ_2[x]/J$                                    & $1 + J$         \\
    $x+1$                 & $\{0+J, (1+x) + J, (x+x^2) + J, (1 + x^2) +J\}$ & $(1+x) + J$       \\
    $x^2 + x + 1$         & $\{0+J, (1 + x + x^2) + J\}$                    & $(1 + x + x^2) + J$\\
    $x^3+1$               & $\{0 + J\}$                                     & $0 + J$ \\
    \end{tabular}
  \]
\end{proof}

\begin{exercise}
  \label{ex:prime_ideal_of_Z[x]}
  Να δείξετε ότι το ιδεώδες $(3, x^2+1)$ του $\ZZ[x]$ είναι ριζικό.
\end{exercise}

\begin{proof}[Λύση]
  Από την Άσκηση \ref{ex:DF_7.4.31}, αρκεί να δείξουμε ότι το $(3, x^2+1) \in \mSpec(\ZZ[x])$ (γιατί;). Όμοια με την Άσκηση \ref{ex:DF_9.2.7}, παρατηρούμε ότι 
  \[
  \ZZ[x]/(3,x^2+1) \cong \ZZ_3[x]/(x^2+1).
  \]
  Το πολυώυνμο $x^2+1$ είναι ανάγωγο στο $\ZZ_3[x]$, διότι δεν έχει ρίζα στο $\ZZ_3$ και ο βαθμός του είναι δύο\footnote{Δεν είναι δύσκολο να δείξει κανείς ότι ένα πολυώυνο βαθμού δύο ή τρία πάνω με συντελεστές πάνω από ένα σώμα είναι ανάγωγο αν και μόνο αν δεν έχει ρίζα στο σώμα αυτό (βλ., για παράδειγμα, \cite[Πρόταση 9.4.9]{DF04}).}. Επομένως, από το Λήμμα 2.14 (2) έπεται ότι $(x^2+1) \in \mSpec(\ZZ_2[x])$ και γι αυτό, από την Πρόταση 1.10, ο $\ZZ_3[x]/(x^2+1)$ είναι σώμα. Κατά συνέπεια, ο $\ZZ[x]/(3,x^2+1)$ είναι σώμα και γι αυτό $(3, x^2+1) \in \mSpec(\ZZ[x])$, πάλι από την Πρόταση 1.10, το οποίο ήταν το ζητούμενο.
\end{proof}

\begin{remark}
  Η λύση της Άσκησης \ref{ex:prime_ideal_of_Z[x]} δουλεύει για κάθε ιδεώδες της μορφής $(p, f(x))$, όπου $p$ είναι πρώτος αριθμός και $f(x) \in \ZZ[x]$ είναι τέτοιο ώστε η αναγωγή του modulo $p$ είναι ανάγωγο στοιχείο του $\ZZ_p[x]$. 

  Επίσης, από το Θεώρημα 3.6, ο $\ZZ[x]$ είναι περιοχή μονοσήμαντης παραγοντοποίησης και κατά συνέπεια $(f(x)) \in \Spec(\ZZ[x])$, για κάθε ανάγωγο πολυώνυμο $f(x) \in \ZZ[x]$, από το Λήμμα 2.14. Τέλος, $(0) \in \Spec(\ZZ[x])$, διότι ο $\ZZ[x]$ είναι ακέραια περιοχή.

  Αποδεινύεται ότι αυτά είναι \emph{όλα} τα στοιχεία του $\Spec(\ZZ[x])$ (βλ., για παράδειγμα, \cite[Πρόταση 1.4.1]{Mal08}). Κάτι αντίστοιχο ισχύει και για τον πολυωνυμικό δακτύλιο $F[x,y]$, όπου $F$ είναι σώμα (βλ., για παράδειγμα, \cite[Πρόταση 1.4.2]{Mal08}).
\end{remark}

\begin{thebibliography}{99}
    \bibitem{DF04}
    D. S. Dummit και R. M. Foote, 
    \textit{Abstract algebra},
    third edition, Wiley, 2004. 
    \bibitem{Mal08}
    Μ. Μαλιάκας, 
    \textit{Εισαγωγή στην Μεταθετική Άλγεβρα}, 
    Εκδόσεις Σοφία, 2008.
\end{thebibliography}
\end{document}